\section*{第 \S\arabic{exercisegroup} 节习题}
\phantomsection
\addcontentsline{toc}{section}{第 \protect\S\arabic{exercisegroup} 节习题}

\begin{enumerate}
\def\labelenumi{\arabic{enumi})}
\tightlist
\item
  \begin{enumerate}
  \def\labelenumii{\alph{enumii})}
  \tightlist
  \item
    设 H 是紧区间 \([a, b] \subset R\) 上的凸函数所成之集；假设集 \(\mathrm{H}(a)\) 与 \(\mathrm{H}(b)\) 在 R 中有上界，并且存在一点 c，使得 a \textless{} c \textless{} b 且 \(\mathrm{H}(c)\) 在 R 中有下界；证明 H 是 {]}a, b{[} 上的等度连续集（Gen.~Top., X, p.~283）。
  \end{enumerate}
\end{enumerate}

\begin{enumerate}
\def\labelenumi{\alph{enumi})}
\setcounter{enumi}{1}
\tightlist
\item
  设 H 是区间 \(I \subset R\) 上的凸函数所成之集，F 是 H 上的一个滤子，它在 I 上逐点收敛于函数 \(f_{0}\)；证明 F 在含于 I 的每个紧区间上一致收敛于 \(f_{0}\)。
\end{enumerate}

\begin{enumerate}
\def\labelenumi{\arabic{enumi})}
\setcounter{enumi}{1}
\item
  证明每个凸函数 \(f\)（在紧区间 \(\mathrm{I} \subset \mathbf{R}\) 上）都是 \(\mathrm{I}\) 上的一个递减的一致收敛凸函数序列的极限，这些函数在 \(\mathrm{I}\) 上具有二阶导数（先考虑函数 \((x - a)^+\)，并把 \(f\) 用一个仿射线性函数与一个线性组合 \(\sum c_j(x - a_j)^+\)（系数为 \(c_j \geqslant 0\)）之和来逼近）。
\item
  设 \(f\) 是区间 \(I \subset \mathbf{R}\) 上的凸函数。
\end{enumerate}

\begin{enumerate}
\def\labelenumi{\alph{enumi})}
\item
  证明若 \(f\) 不是常数，它就不能在 \(I\) 的内点处达到其上确界。
\item
  证明若 I 在 R 中相对紧，则 f 在 I 上有下界。
\item
  证明若 \(\mathbf{I} = \mathbf{R}\) 且 \(f\) 不是常数，则 \(f\) 在 \(\mathbf{I}\) 上没有上界。
\end{enumerate}

\begin{enumerate}
\def\labelenumi{\arabic{enumi})}
\setcounter{enumi}{3}
\item
  为使函数 \(f\) 在紧区间 \([a, b] \subset \mathbf{R}\) 上凸，必须且只需它在 \(]a, b[\) 上凸，并且有 \(f(a) \geqslant f(a+)\) 与 \(f(b) \geqslant f(b-)\)。
\item
  设 \(f\) 是开区间 \(]a, +\infty[\) 上的凸函数；若存在一点 \(c > a\)，使得 \(f\) 在 \(]c, +\infty[\) 上严格递增，则 \(\lim_{x \to +\infty} f(x) = +\infty\)。
\item
  设 \(f\) 是区间 \(]a, +\infty[\) 上的凸函数；证明 \(f(x)/x\) 具有一个极限（有限或等于 \(+\infty\)），当 \(x\) 趋于 \(+\infty\) 时；这个极限也是 \(f_d'(x)\) 与 \(f_g'(x)\) 的极限；则它是 \(>0\)，如果 \(f(x)\) 趋于 \(+\infty\)（当 \(x\) 趋于 \(+\infty\) 时）。
\item
  设 \(f\) 是区间 \(]a, b[\) 上的凸函数，其中 \(a \geqslant 0\)；证明在此区间上函数 \(x \mapsto f(x) - xf'(x)\)（即在点 \(x\) 处与 \(f\) 的图像相切的右半切线的“原点纵坐标”）是递减的（若 \(f\) 严格凸则严格递减）。
\end{enumerate}

由此推出：

\(a)\) 若 \(f\) 在点 \(a\) 处具有有限的右极限，则 \((x - a)f_d'(x)\) 在该点处有等于 0 的右极限。

\begin{enumerate}
\def\labelenumi{\alph{enumi})}
\setcounter{enumi}{1}
\item
  在 {]}a, b{[} 上，或者 \(f(x)/x\) 递增，或者 \(f(x)/x\) 递减，或者存在 \(c \in ]a\) , b{[}，使得 \(f(x)/x\) 在 {]}a, c{[} 上递减而在 {]}c, b{[} 上递增。
\item
  假设 \(b = +\infty\)：证明若
\end{enumerate}

\[
\beta = \lim _ {x \rightarrow + \infty} \left(f (x) - x f _ {d} ^ {\prime} (x)\right)
\]

有限，则 \(\alpha = \lim_{x \to +\infty} f(x) / x\) 也有限，并且直线 \(y = \alpha x + \beta\) 渐近 \(^{5}\) 于 f 的图像，且位于该图像之下（若 f 严格凸则严格位于其下）。

\begin{enumerate}
\def\labelenumi{\arabic{enumi})}
\setcounter{enumi}{7}
\tightlist
\item
  设 \(f\) 是开区间 \(I \subset \mathbf{R}\) 上的上半连续实值有限函数。这时 \(f\) 凸当且仅当 \(\limsup_{h \to 0, h \neq 0} \frac{f(x + h) + f(x - h) - 2f(x)}{h^2} \geqslant 0\) 对一切 \(x \in I\) 成立。（先证明：对一切 \(\varepsilon > 0\)，函数 \(f(x) + \varepsilon x^2\) 是凸的，使用 I, p.~31 的命题 9。）
\end{enumerate}

¶ 9) 设 \(f\) 是区间 \(\mathrm{I} \subset \mathbf{R}\) 上的下半连续实值有限函数。为使 \(f\) 在 \(\mathrm{I}\) 上凸，只需：对每一对点 \(a, b\)（属于 \(\mathrm{I}\)，满足 \(a < b\)），存在一点 \(z\)，使得 \(a < z < b\)，并且 \(\mathbf{M}_z\) 位于线段 \(\mathbf{M}_a\mathbf{M}_b\) 之下（用反证法论证，注意由点 \(x\)（使 \(\mathbf{M}_x\) 严格位于 \(\mathbf{M}_a\mathbf{M}_b\) 之上者）所成的集是开集）。

¶ 10) 设 \(f\) 是定义在区间 \(\mathrm{I}\subset \mathbf{R}\) 上的实值有限函数，使得

\[
f \left(\frac {x + y}{2}\right) \leqslant \frac {1}{2} (f (x) + f (y))
\]

对 I 中一切 x, y 成立。证明若 f 在含于 I 的某个开区间 {]}a, b{[} 上有上界，则 f 在 I 上凸（先证明 f 在含于 I 的每个紧区间上有上界，再证明 f 在 I 的每个内点处连续）。

¶ 11) 设 \(f\) 是开区间 \(\mathrm{I} \subset \mathbf{R}\) 上的连续函数，在 \(\mathrm{I}\) 的每一点都有有限右导数。若对每个 \(x \in \mathrm{I}\) 与每个 \(y \in \mathrm{I}\)（满足 \(y > x\)），点 \(\mathbf{M}_y = (y, f(y))\) 都位于 \(f\) 的图像在点 \(\mathbf{M}_x = (x, f(x))\) 处的右半切线之上，证明 \(f\) 在 \(\mathrm{I}\) 上凸（用中值定理证明 \(f(y) = f(x)\)）。

\[
f _ {d} ^ {\prime} (y) \geqslant \frac {f (y) - f (x)}{y - x} \text {   for   } x <   y).
\]

给出一个函数的例子，它不凸、处处有有限右导数，并且对每个 \(x \in I\) 存在依赖于 x 的数 \(h_{x} > 0\)，使得 \(M_{y}\) 位于点 \(M_{x}\) 处的右半切线之上（对一切满足 \(x \leqslant y \leqslant x + h_{x}\) 的 y）。然而若进一步假设 f 在 I 上可导，则后一条件对 f 的凸性是充分的（利用 I, p.~12 的推论）。

¶ 12) 设 \(f\) 是开区间 \(\mathrm{I}\subset \mathbf{R}\) 上的连续实函数；假设对 I 的每一对点 \((a,b)\)（满足 \(a < b\)），\(f\) 的图像或者整个位于线段 \(\mathbf{M}_a\mathbf{M}_b\) 之上，或者整个位于其下（在区间 \([a,b]\) 上）。证明 \(f\) 在整个 I 上凸，或者在整个 I 上凹（若在 \(]a\)、\(b[\) 中存在一点 \(c\)，使得 \(\mathbf{M}_c\) 严格位于线段 \(\mathbf{M}_a\mathbf{M}_b\) 之上，则证明对每个 \(x\in \mathrm{I}\)（满足 \(x > a\)），\(f\) 的图像位于线段 \(\mathbf{M}_a\mathbf{M}_x\) 之上（在区间 \([a,x]\) 上））。

\begin{enumerate}
\def\labelenumi{\arabic{enumi})}
\setcounter{enumi}{12}
\item
  设 \(f\) 是开区间 \(\mathrm{I} \subset \mathbf{R}\) 上的可导实函数。假设对每一对点 \((a, b)\)（属于 \(\mathrm{I}\) 且满足 \(a < b\)），存在唯一点 \(c \in ]a, b[\)，使得 \(f(b) - f(a) = (b - a)f'(c)\)；证明 \(f\) 在 \(\mathrm{I}\) 上严格凸或在 \(\mathrm{I}\) 上严格凹（证明 \(f'\) 在 \(\mathrm{I}\) 上严格单调）。
\item
  设 \(f\) 是开区间 \(I \subset \mathbf{R}\) 上的凸实函数且严格单调；设 \(g\) 是 \(f\) 的反函数（定义在区间 \(f(I)\) 上）。证明若 \(f\) 在 \(I\) 上递减（相应地，递增），则 \(g\) 在 \(f(I)\) 上凸（相应地，凹）。
\item
  设 I 是含于 \(]0, +\infty[\) 的区间；证明若 \(f(1 / x)\) 在 I 上凸，则 \(xf(x)\) 亦然，反之亦然。
\end{enumerate}

* 16) 设 \(f\) 是 \(]0, +\infty[\) 上的正凸函数，而 \(a, b\) 是两个任意实数。证明在下列情形中函数 \(x^a f(x^{-b})\) 在 \(]0, +\infty[\) 上凸：

\[
1 ^ {\circ} a = \frac {1}{2} (b + 1), | b | \geqslant 1;
\]

\(2^{\circ} x^{a} f(x^{-b})\) 递增，\(a(b - a) \geqslant 0\)，\(a \geqslant \frac{1}{2}(b + 1)\)；

\(3^{\circ} x^{a} f(x^{-b})\) 递减，\(a(b-a) \geqslant 0\)，\(a \leqslant \frac{1}{2}(b+1)\)。

在对 \(f\) 的同样假设下，证明 \(e^{x / 2}f(e^{-x})\) 是凸的（利用 I, p.~45 的习题 2）。*

\begin{enumerate}
\def\labelenumi{\arabic{enumi})}
\setcounter{enumi}{16}
\item
  设 \(f\) 与 \(g\) 是区间 \(I = [a, b]\) 上的两个正凸函数；假设存在数 \(c \in \mathrm{I}\)，使得在区间 \([a, c]\) 与 \([c, b]\) 的每一个中，函数 \(f\) 与 \(g\) 同向变化。证明乘积 \(fg\) 在 \(\mathrm{I}\) 上凸。
\item
  设 \(f\) 是区间 \(\mathrm{I} \subset \mathbf{R}\) 上的凸函数，\(g\) 是包含 \(f(\mathrm{I})\) 的区间上的递增凸函数；证明 \(g \circ f\) 在 \(\mathrm{I}\) 上凸。
\end{enumerate}

¶ 19) 设 \(f\) 与 \(g\) 是两个实值有限函数，\(f\) 在区间 I 上定义且连续，而 \(g\) 在 \(\mathbf{R}\) 上定义且连续。假设对每对实数 \((\lambda, \mu)\)，函数 \(g(f(x) + \lambda x + \mu)\) 都在 I 上凸。

\begin{enumerate}
\def\labelenumi{\alph{enumi})}
\item
  证明 \(g\) 在 \(\mathbf{R}\) 上凸且单调。
\item
  若 g 在 R 上递增（相应地，递减），证明 f 在 I 上凸（相应地，凹）（利用命题 7）。
\end{enumerate}

\begin{enumerate}
\def\labelenumi{\arabic{enumi})}
\setcounter{enumi}{19}
\item
  证明集 \(\mathfrak{K}\)（区间 \(\mathrm{I} \neq \mathbf{R}\) 上的凸函数所成）对于“\(f(x) \leqslant g(x)\)”这一序（对一切 \(x \in \mathrm{I}"\)）构成一个格（Set Theory, III, p.146）。给出两个凸函数 \(f\)、\(g\)（在 \(\mathrm{I}\) 上）的例子，使它们的下确界在 \(\mathfrak{K}\) 中在某些点取与 \(\inf(f(x), g(x))\) 不同的值。给出一个无限族 \((f_{\alpha})\)（由 \(\mathfrak{K}\) 中的函数组成）的例子，使得 \(\inf_{\alpha} f_{\alpha}(x)\) 在每一点 \(x \in \mathrm{I}\) 都有限，然而 \(\mathfrak{K}\) 中没有函数小于所有 \(f_{\alpha}\)。
\item
  设 \(f\) 是开区间 \(\mathbf{I} \subset \mathbf{R}\) 上的上半连续实值有限函数。为使 \(f\) 在 \(\mathbf{I}\) 上严格凸，必须且只需不存在直线局部位于 \(\mathbf{G}\)（\(f\) 的图像）上方（在 \(\mathbf{G}\) 的一点处）。
\item
  设 \(f_{1}, \ldots, f_{n}\) 是紧区间 \(I \subset \mathbf{R}\) 上的连续凸函数；假设对一切 \(x \in I\) 有 \(\sup(f_{j}(x)) \geqslant 0\)。证明存在 \(n\) 个数 \(\alpha_{j} \geqslant 0\)，使得 \(\sum_{j=1}^{n} \alpha_{j} = 1\)，并且 \(\sum_{j=1}^{n} \alpha_{j} f_{j}(x) \geqslant 0\) 在 \(I\) 上成立。（先处理 \(n = 2\) 的情形，考虑一个
\end{enumerate}

点 \(x_{0}\)（上包络 \(\sup(f_{1}, f_{2})\) 在其处达到其极小值）；当 \(x_{0}\) 是 I 的内点时，确定 \(\alpha_{1}\)，使 \(\alpha_{1}f_{1} + (1 - \alpha_{1})f_{2}\) 的左导数在 \(x_{0}\) 处为零。通过对 n 的归纳过渡到一般情形；对 \(f_{1}, \ldots, f_{n-1}\) 在使 \(f_{n}(x) \leqslant 0\) 的紧区间上的限制使用归纳假设。

\begin{enumerate}
\def\labelenumi{\arabic{enumi})}
\setcounter{enumi}{22}
\item
  设 \(f\) 是紧区间 \(\mathrm{I} \subset \mathbf{R}\) 上的连续实函数；在函数 \(g \leqslant f\)（在 \(\mathrm{I}\) 上凸者）之中存在一个 \(g_0\)，它大于所有其他的函数。设 \(\mathrm{F} \subset \mathrm{I}\) 是 \(x \in \mathrm{I}\)（在其处 \(g_0(x) = f(x)\)）所成之集；证明 \(\mathrm{F}\) 非空，并且在与 \(\mathrm{F}\) 邻接的每个开区间上函数 \(g_0\) 等于一个仿射线性函数（用反证法论证）。
\item
  设 \(\mathrm{P}(x)\) 是 \(n\) 次多项式，实系数，其全部根都是实的且含于区间 \([-1, 1]\)。设 \(k\) 是使 \(1 \leqslant k \leqslant n\) 的整数。证明有理函数
\end{enumerate}

\[
f (x) = x + \frac {\mathrm{P} ^ {(k - 1)} (x)}{\mathrm{P} ^ {(k)} (x)}
\]

在 R 中其有定义的每个区间上递增；若 \(c_{1} < c_{2} < \cdots < c_{r}\) 是它的极点（含于 \([-1, 1]\)），则 f 对 \(x < c_{1}\) 凸而对 \(x > c_{r}\) 凹。由此推出：当 a 跑遍 \([-1, 1]\) 时，把 \(k^{th}\) 阶导数（\((x - a)\mathrm{P}(x)\) 的）的零点全部包含在内的最大区间的长度在 a = 1 或 a = -1 时达到其最大值。

\begin{enumerate}
\def\labelenumi{\arabic{enumi})}
\setcounter{enumi}{24}
\tightlist
\item
  我们称实函数 \(f\)（定义在 \([0, +\infty[\) 上）是超加性的，如果 \(f(x + y) \geqslant f(x) + f(y)\) 对 \(x \geqslant 0, y \geqslant 0\) 成立，并且 \(f(0) = 0\)。
\end{enumerate}

\begin{enumerate}
\def\labelenumi{\alph{enumi})}
\item
  给出不连续的超加性函数的例子。
\item
  证明每个凸函数 \(f\)（在 \([0, +\infty[\) 上且满足 \(f(0) = 0\)）都是超加性的。
\item
  若 \(f_{1}\) 与 \(f_{2}\) 都是超加性的，则 \(\inf (f_1, f_2)\) 亦然；利用这一点，给出非凸的连续超加性函数的例子。
\item
  若 f 连续且 \(\geqslant0\) 于区间 \([0,a]\)（a\textgreater0）上，使得 \(f(0)=0\) 与 \(f(x/n)\leqslant f(x)/n\) 对每个整数 \(n\geqslant1\) 成立，证明 f 在点 0 处有右导数（用反证法论证）。特别地，每个满足 \(\geqslant0\) 的连续超加性函数都在点 0 处具有右导数。
\end{enumerate}

